\documentclass{article}
\usepackage[T1]{fontenc}
\usepackage{lmodern}
\usepackage{graphicx}
\usepackage{amsmath,amssymb}
\usepackage[a4paper,margin=2cm]{geometry}
\usepackage[absolute,overlay]{textpos}
\setlength{\TPHorizModule}{1mm}
\setlength{\TPVertModule}{1mm}
\usepackage[indonesian]{babel}
\usepackage{hyperref}
\setlength{\emergencystretch}{2em}
% unit: Halaman 12

\begin{document}

% === Halaman 12 (llm) ===

\section*{Cipher Sederhana dengan operasi XOR}

\begin{itemize}
    \item Sama seperti \textit{Vigenere Cipher}, tetapi dalam mode bit
    \item Setiap bit plainteks di-XOR-kan dengan setiap bit kunci.
\end{itemize}

\[ \text{Enkripsi: } C = P \oplus K \]
\[ \text{Dekripsi: } P = C \oplus K \]

\vspace{5mm}

Contoh:
\begin{tabular}{lll}
plainteks & 01100101 & (karakter 'e') \\
kunci & 00110101 $\oplus$ & (karakter '5') \\
\hline
cipherteks & 01010000 & (karakter 'P') \\
kunci & 00110101 $\oplus$ & (karakter '5') \\
\hline
plainteks & 01100101 & (karakter 'e')
\end{tabular}

\end{document}
